% IEEE conference A4 (IEEEtran). Compile with: pdflatex conference_paper.tex
% Content matches NeuroTrace-DAG_IEEE_A4_Conference.docx
% Confirmatory H-rev: pending until OSF + --touch-test; fill from JSON only.
\documentclass[conference,a4paper]{IEEEtran}
\IEEEoverridecommandlockouts
\usepackage{cite}
\usepackage{amsmath,amssymb,amsfonts}
\usepackage{booktabs}
\usepackage{graphicx}
\usepackage{textcomp}
\usepackage{xcolor}
\usepackage{url}
\usepackage{hyperref}

\begin{document}

\title{When Next-Item Metrics Mislead: Evaluation Validity Across Educational Interaction Logs}

\author{\IEEEauthorblockN{Oscar Man Shrestha}
\IEEEauthorblockA{\textit{Dept.\ name of organization (Affiliation)}\\
City, Country\\
email address or ORCID}
}

\maketitle

\begin{abstract}
Recall@K is the usual report card for next-item models on student logs.
We are less sure it measures what people think it measures.
On Junyi Academy, a gated ranker gets Recall@5 $0.891$ $[0.880, 0.902]$, yet Recent-5 is already $0.885$---within a $0.01$ practical gap---so we do not claim a recommender win.
The more interesting question is whether the \emph{action unit} (question, skill, or topic cluster) changes which simple model looks better.
We lock a small confirmatory comparison of a GRU probe against first-order Markov on Junyi timed, ASSISTments~2009, and XES3G5M, at native units and at matched clusters ($\sim$120), under a freeze holdout and Bonferroni family $m_{\mathrm{planned}}=6$.
Holdout cells stay blank until external registration and one touch-test; Junyi ranking numbers below are exploratory.
\end{abstract}

\begin{IEEEkeywords}
educational data mining, evaluation validity, next-item recommendation,
knowledge tracing logs, preregistration, action granularity
\end{IEEEkeywords}

\section{Introduction}
On educational logs, next-item Recall@K often rewards repeats and the curriculum's natural order.
That is not the same as measuring whether a model captured learning structure.
Worse, changing how we define an ``item''---question id versus skill or topic cluster---can flip which baseline wins.

Our goal here is modest: check whether that unit choice systematically reverses a GRU-versus-Markov contrast under a locked protocol.
We are not proposing a new recommender champion.

\noindent\textbf{Registered question (prereg-v1).}
Does the sign of the paired GRU$-$Markov non-repeat Recall@5 contrast flip between native and cluster units on at least one dataset?
Until registration and a single freeze holdout run, those cells remain pending.

\section{Related Work}
KT and educational recommender papers usually score platform-native ids~\cite{hidasi2016,kang2018}.
Duplicates, unit mismatch, and curriculum order show up often as caveats; we treat them as the object of study.
Library-faithful SASRec/GRU4Rec are outside this confirmatory packet for now; the probes below are diagnostic.

\section{Datasets and Units}
Junyi: 835 concepts, 978 edges, 27{,}434/6{,}290 train/test sequences; embedding pair-AUC vs random 0.904.
We also use ASSISTments skills and XES questions, plus registered maps that collapse items into roughly 120 clusters (two-method intersection on Junyi/ASSISTments; co-occurrence only on XES).

\begin{table}[t]
\centering
\caption{Registered action units.}
\label{tab:units}
\begin{tabular}{@{}lll@{}}
\toprule
Dataset & Native & Cluster \\
\midrule
Junyi timed & topic/exercise & text $k$-means (+cooc) \\
ASSISTments & skill token & text (+cooc) \\
XES3G5M & question & cooc only \\
\bottomrule
\end{tabular}
\end{table}

\section{Methods}
\subsection{Confirmatory comparison}
Freeze seed 20261004; GRU seeds 20261101/02/03; frozen config and fixed \texttt{best\_epoch}; $\ge$10{,}000 sequence bootstrap; $m_{\mathrm{planned}}=6$ Bonferroni; $\ge$500 non-repeat queries per cell.
A dataset supports the primary claim only if native and cluster both show \texttt{supported\_*} verdicts with opposite signs; the project needs $\ge$1 such dataset.
Models in the confirmatory tables: popularity, recency, \texttt{markov\_order1}, \texttt{gru\_probe}.

\subsection{Exploratory Junyi work}
Embeddings, memory classifiers, the DAG, the review/advance gate, and the gated ranker are useful diagnostics.
They are not part of the confirmatory family.

\section{Results}
\subsection{Primary confirmatory cells (pending)}
Tag \texttt{prereg-v1} freezes the protocol and artifacts.
Table~\ref{tab:hrev} is empty on purpose until OSF/Zenodo and one \texttt{--touch-test}.
Numbers will come only from \texttt{primary\_claims} JSON.

\begin{table}[t]
\centering
\caption{Confirmatory unit cells (pending touch-test).}
\label{tab:hrev}
\begin{tabular}{@{}llcc@{}}
\toprule
Dataset & Unit & Verdict & CI \\
\midrule
junyi\_timed & native/cluster & pending & --- \\
assistments & native/cluster & pending & --- \\
xes3g5m & native/cluster & pending & --- \\
\bottomrule
\end{tabular}
\end{table}

\subsection{Exploratory Junyi ranking}
Capped RAGR R@5 $=0.891$ $[0.880, 0.902]$ vs Recent-5 $0.885$ $[0.872, 0.897]$ ($n{=}800$).
Full-test RAGR R@5 $=0.889$ $[0.885, 0.893]$ ($n{=}6290$).
Revisit logistic AUC $=0.864$; gate Brier $=0.0928$, ECE $=0.0335$.
The sub-$0.01$ gap is why we do not sell this as a ranking breakthrough.

\subsection{Cross-dataset diagnostics}
Markov looks strong on curriculum-heavy XES advance slices; GRU probes move around by dataset.
That pattern is why we preregistered the sign-flip test; it is not itself the confirmatory answer.

\section{Discussion and Limitations}
A null or mixed result is still a result: it would mean the preregistered flip did not show up under the locked rules.
Missing library SASRec, unverified redistribution licences, and disagreeing cluster maps (low ARI) are real limits.

\section{Conclusion}
We locked a small confirmatory test about action units and left the Junyi product numbers clearly exploratory.
After registration, one touch-test fills the holdout table; until then those cells stay blank.

\section*{Acknowledgment}
Code and prereg packet: \url{https://github.com/Oscar-man-shrestha/eval-validity-edu} (tag \texttt{prereg-v1}).

\begin{thebibliography}{00}
\bibitem{hidasi2016} B. Hidasi et al., ``Session-based recommendations with recurrent neural networks,'' in \emph{Proc. ICLR}, 2016.
\bibitem{kang2018} W.-C. Kang and J. McAuley, ``Self-attentive sequential recommendation,'' in \emph{Proc. IEEE ICDM}, 2018, pp. 197--206.
\bibitem{edudata} F. Wang et al., ``EduData,'' GitHub, accessed 2026.
\bibitem{osf} Center for Open Science, ``OSF Preregistration,'' \url{https://osf.io/registries}, 2026.
\bibitem{preregv1} NeuroTrace-DAG, ``prereg-v1 release,'' 2026. [Online]. Available: \url{https://github.com/Oscar-man-shrestha/eval-validity-edu/releases/tag/prereg-v1}
\end{thebibliography}

\end{document}
